\documentclass[11pt,a4paper]{article}
\usepackage{DDTA_DERMATRIAGE_DOCUMENTATION_STYLE_R6}
\begin{document}
\selectlanguage{italian}
\renewcommand{\candidateauthority}{\texttt{CURRENT\_GOVERNED}}
\hypersetup{pdftitle={DermaTriage Governed DDTA Documentation - R1}}
\fancyhead[R]{\sffamily\small\color{DDTAGray}DERMATRIAGE-GOV-R1}

\begin{titlepage}
\thispagestyle{empty}
\vspace*{9mm}
{\sffamily\bfseries\color{DDTANavy}\fontsize{15}{18}\selectfont Documentation-Driven Threat Analysis}\par
\vspace{5mm}
{\sffamily\bfseries\color{DDTANavy}\fontsize{28}{32}\selectfont DermaTriage Governed DDTA Documentation}\par
\vspace{2mm}
{\sffamily\color{DDTABlue}\fontsize{18}{22}\selectfont Governed Baseline R1 - Native DDTA Project Documentation}\par

\vspace{10mm}
\begin{tcolorbox}[enhanced,arc=3pt,boxrule=0pt,colback=DDTANavy,left=12pt,right=12pt,top=10pt,bottom=10pt]
{\color{white}\sffamily\bfseries Current governed DDTA project documentation}\par
\vspace{2mm}
{\color{white!88}\small Gerarchia leggibile in profondità: MacroRequirement $\rightarrow$ Decision $\rightarrow$ FunctionalRequirement $\rightarrow$ eventuale specializzazione. Componenti, processi, algoritmi, sistemi e tecnologie materialmente rilevanti sono espressi mediante gli stessi elementi DDTA pertinenti, senza campi tecnici paralleli.}
\end{tcolorbox}

\vspace{8mm}
\begin{tabularx}{\textwidth}{L{46mm}Y}
\toprule
\textbf{Status} & \texttt{CURRENT\_GOVERNED} / PRIMARY BA SOURCE ALLOWED \\
\textbf{Coverage} & Project framing + governed MR/Decision/FR baseline with explicit technical/project choices \\
\textbf{Methodology authority} & \texttt{DDTA\_DOCUMENTATION\_BA\_AUTHORING\_GUIDE\_R5} \\
\textbf{Promotion predecessor} & \texttt{bf2236c} \\
\textbf{Authoritative source package} & \texttt{DermaTriage-Docs-20260830T152637Z-1-001.zip} \\
\textbf{Source package SHA-256} & {\scriptsize\nolinkurl{E9ED2C507BEFB95F54A52084687CD1E8798863AE81CF69D09568864D8CBF280E}} \\
\bottomrule
\end{tabularx}

\vfill
{\sffamily\small\color{DDTAGray}5 settembre 2026}
\end{titlepage}

\section*{Indice gerarchico del progetto}
\addcontentsline{toc}{section}{Indice gerarchico del progetto}

\begin{hierarchybox}[title={DermaTriage - current governed hierarchy}]
{\small\textbf{Legenda:} MR = MacroRequirement; DEC = Decision; FR = FunctionalRequirement; STOP = decomposizione intenzionalmente arrestata; SUPERSEDED = identità storica non attiva.}
\begin{lstlisting}[basicstyle=\ttfamily\scriptsize]
DERMATRiAGE
|
|-- MR-01  Valutazione di triage del caso dermatologico
|   |
|   |-- DEC-12  Pipeline AI image-based a quattro stadi
|   |   |-- FR-16  Classificazione image-based con EfficientNet-B4
|   |   |-- FR-17  Descrizione clinica con Qwen2-VL-7B-Instruct
|   |   |-- FR-18  Retrieval storico con ChromaDB/all-MiniLM-L6-v2
|   |   `-- FR-19  Sintesi multi-source con BioMistral-7B
|   |
|   |-- DEC-13  Training della baseline EfficientNet-B4
|   |   `-- FR-20  Addestramento della baseline classificatoria
|   |
|   |-- DEC-14  Esposizione del servizio e integrazione B4
|   |   |-- FR-21  Analisi diretta tramite POST /analyze
|   |   `-- FR-22  Diagnosi integrata B4 tramite POST /diagnose
|   |
|   |-- DEC-01  Continuita del triage in assenza di immagine
|   |   `-- FR-01  Determinazione urgenza su base sintomatologica
|   |
|   `-- DEC-02  Adozione della P-scale per la priorita operativa
|       `-- FR-02  Derivazione della priorita P-scale
|
|-- MR-02  Indirizzamento specialistico
|   `-- STOP AT MR
|
|-- MR-03  Gestione della validazione clinica degli esiti
|   |-- DEC-03  Separazione tra esito originario e revisione clinica
|   |   |-- FR-03  Registrazione correlata della revisione clinica
|   |   `-- FR-12  Distinzione tra conferma e correzione
|   |-- DEC-15  Scambio di validazione con B4
|   |   `-- FR-23  Write-back/retrieval della validazione B4
|   `-- DEC-16  Autenticazione delle interfacce di review
|       |-- FR-24  X-API-Key per operazioni amministrative DermaTriage
|       `-- FR-25  Bearer JWT per interazione B4
|
`-- MR-04  Adattamento controllato sulla base della revisione clinica
    |     dependsOn MR-03
    |-- DEC-04 -> FR-04, FR-05
    |-- DEC-05 -> FR-13, FR-14, FR-15 [FR-11 SUPERSEDED]
    |-- DEC-06 -> FR-09
    |-- DEC-07 -> no active concrete SpecializedRequirement
    |-- DEC-08 -> FR-10
    |-- DEC-09 -> FR-06
    |-- DEC-10 -> FR-07
    |-- DEC-11 -> FR-08
    |-- DEC-17  Gestione prompt con PromptManager
    |   `-- FR-26  Versioning e persistenza dei prompt
    `-- DEC-18  Retraining incrementale EfficientNet-B4
        `-- FR-27  Fine-tuning degli ultimi blocchi e classifier head
\end{lstlisting}
\end{hierarchybox}

\begin{statusbox}[title={Reading convention}]
Gli ID sono identità stabili e non impongono un ordine di lettura. Il corpo del documento segue intenzionalmente una proiezione \textbf{depth-first}: per ogni MR vengono presentate una Decision e tutti i suoi FR prima di passare alla Decision successiva; completato il branch si passa al successivo MR.
\end{statusbox}

\tableofcontents
\clearpage

\section{Status, authority e convenzione documentale}

\begin{projectbox}[title={Authority statement}]
Ogni elemento attivo di questa baseline espone \textbf{Lifecycle}, \textbf{Authority}, \textbf{Type} e \textbf{Title}. Il valore \candidateauthority{} indica che l'elemento appartiene alla baseline corrente \texttt{DERMATRIAGE-GOV-R1}, autorizzata come sorgente primaria per la Base Analysis.
\end{projectbox}

\section{Project problem framing}

\begin{projectbox}
DermaTriage affronta il problema di supportare il triage precoce di casi dermatologici relativi a lesioni cutanee potenzialmente oncologiche, utilizzando le informazioni disponibili sul caso per determinare una priorità di presa in carico e supportare l'instradamento verso l'assistenza specialistica appropriata.
\end{projectbox}

\subsection{Scope boundary}

\textbf{Governed in the current baseline:}
\begin{itemize}
\item valutazione di urgenza e priorità orientata al triage;
\item responsabilità macro di indirizzamento specialistico;
\item gestione degli esiti di revisione clinica;
\item adattamento controllato utilizzando evidence derivata dalla revisione clinica.
\end{itemize}

\textbf{Not established as governed project authority:}
\begin{itemize}
\item autorità per diagnosi clinica definitiva;
\item responsabilità completa di specialist selection/booking;
\item deployment automatico o autorità finale di promotion dopo qualification;
\item semantica normativa completa di ogni valore di configurazione o dettaglio di realization; i fatti tecnici source-supported rimangono comunque descritti quando materialmente rilevanti.
\end{itemize}

\section{MR-01 - Valutazione di triage del caso dermatologico}\label{mr-01}

\begin{mrmeta}
\begin{tabularx}{\linewidth}{L{43mm}Y}
\toprule
\textbf{ID} & \texttt{MR-01} \\
\textbf{Lifecycle} & \texttt{current} \\
\textbf{Authority} & \candidateauthority \\
\textbf{Type} & \MR \\
\textbf{Title} & Valutazione di triage del caso dermatologico \\
\bottomrule
\end{tabularx}
\end{mrmeta}

\subsection{Intent}
Determinare, dalle informazioni disponibili sul caso dermatologico, una valutazione di urgenza utile a stabilire la priorità di presa in carico specialistica.

\subsection{Context}
La responsabilità riguarda la valutazione di urgenza e priorità di un caso dermatologico relativo a una lesione cutanea potenzialmente oncologica. Il caso può essere disponibile con o senza immagine; la responsabilità di triage non coincide con l'autorità per una diagnosi clinica definitiva e non possiede l'intero processo di instradamento specialistico.

\subsection{Stakeholders}
\begin{itemize}
\item professionisti sanitari che utilizzano o revisionano gli esiti di triage;
\item responsabili della successiva presa in carico specialistica;
\item soggetti cui il caso dermatologico si riferisce;
\item responsabili del servizio DermaTriage.
\end{itemize}

\subsection{Scope}
\textbf{In scope:} valutazione dell'urgenza del caso; produzione delle informazioni analitiche necessarie al triage; rappresentazione della priorità operativa P1-P4; supporto del percorso image-based e del fallback senza immagine.

\textbf{Out of scope:} diagnosi clinica definitiva; scelta completa dello specialista o prenotazione; gestione della presa in carico specialistica; semantica completa degli SLA finché non sufficientemente governata.

\subsection{Assumptions / Constraints}
L'immagine non è una precondizione universale del triage. Non viene assunto un vocabolario chiuso universale dei sintomi e DermaTriage non acquisisce autorità diagnostica definitiva.

\subsection{dependsOn}
Nessuna dipendenza macro canonica è stabilita per \texttt{MR-01}.

\subsection{DEC-12 - Adozione della pipeline AI image-based a quattro stadi}\label{dec-12}
\begin{decmeta}
\begin{tabularx}{\linewidth}{L{43mm}Y}
\toprule
\textbf{ID} & \texttt{DEC-12} \\
\textbf{Lifecycle} & \texttt{current} \\
\textbf{Authority} & \candidateauthority \\
\textbf{Type} & \DEC \\
\textbf{Title} & Adozione della pipeline AI image-based a quattro stadi \\
\textbf{Parent MacroRequirement} & \texttt{MR-01} \\
\textbf{Decision kind} & \texttt{SELECTABLE} \\
\bottomrule
\end{tabularx}
\end{decmeta}

\subsubsection{Context}
Quando una \texttt{SkinLesionImage} è disponibile, DermaTriage deve trasformare l'immagine e le altre informazioni del caso in un esito analitico utilizzabile dal successivo mapping operativo di priorità.

\subsubsection{Decision}
DermaTriage adotta una pipeline sequenziale a quattro stadi: \textbf{EfficientNet-B4} esegue la classificazione image-based di urgenza; \textbf{Qwen2-VL-7B-Instruct} produce una descrizione clinica testuale della lesione; \textbf{ChromaDB} con embedding \nolinkurl{sentence-transformers/all-MiniLM-L6-v2}, cosine similarity e retrieval top-5 recupera casi storici simili; \textbf{BioMistral-7B} combina gli output precedenti con le informazioni sintomatologiche per produrre la sintesi AI di triage.

\subsubsection{Consequences}
I quattro stadi hanno ruoli distinti e i loro output devono rimanere correlati allo stesso caso. Il risultato analitico della pipeline resta distinto dalla successiva priorità operativa P1-P4.

\subsubsection{FR-16 - Classificazione image-based di urgenza con EfficientNet-B4}\label{fr-16}
\begin{frmeta}
\begin{tabularx}{\linewidth}{L{43mm}Y}
\toprule
\textbf{ID} & \texttt{FR-16} \\
\textbf{Lifecycle} & \texttt{current} \\
\textbf{Authority} & \candidateauthority \\
\textbf{Type} & \FR \\
\textbf{Title} & Classificazione image-based di urgenza con EfficientNet-B4 \\
\textbf{Parent Decision} & \texttt{DEC-12} \\
\textbf{Owning MR (derived)} & \texttt{MR-01} \\
\textbf{SpecializedRequirement relation} & 0 active instances \\
\bottomrule
\end{tabularx}
\end{frmeta}

\paragraph{Functional obligation.}
Classificare una skin-lesion image nel dominio di urgenza analitica HIGH/MEDIUM/LOW producendo anche la confidence associata.

\begin{normbox}[title={Functional clauses / inherited \texttt{normativeClause [1..*]}}]
Quando una \texttt{SkinLesionImage} è disponibile per il caso, EfficientNet-B4 MUST produrre una \texttt{ImageUrgencyClassification} appartenente a \texttt{HIGH | MEDIUM | LOW} con confidence associata. L'input della baseline è RGB 380x380; la soglia HIGH documentata è \texttt{0.25}.
\end{normbox}

\paragraph{SPO references (L2 support).}
\begin{lstlisting}
SkinLesionImage -- classified by --> EfficientNet-B4
EfficientNet-B4 -- produces --> ImageUrgencyClassification
ImageUrgencyClassification -- includes --> Confidence
\end{lstlisting}

\subsubsection{FR-17 - Produzione della descrizione clinica con Qwen2-VL-7B-Instruct}\label{fr-17}
\begin{frmeta}
\begin{tabularx}{\linewidth}{L{43mm}Y}
\toprule
\textbf{ID} & \texttt{FR-17} \\
\textbf{Lifecycle} & \texttt{current} \\
\textbf{Authority} & \candidateauthority \\
\textbf{Type} & \FR \\
\textbf{Title} & Produzione della descrizione clinica con Qwen2-VL-7B-Instruct \\
\textbf{Parent Decision} & \texttt{DEC-12} \\
\textbf{Owning MR (derived)} & \texttt{MR-01} \\
\textbf{SpecializedRequirement relation} & 0 active instances \\
\bottomrule
\end{tabularx}
\end{frmeta}

\paragraph{Functional obligation.}
Produrre dall'immagine una descrizione clinica testuale utilizzabile dagli stadi successivi della pipeline.

\begin{normbox}[title={Functional clauses / inherited \texttt{normativeClause [1..*]}}]
Per una \texttt{SkinLesionImage} processata dal percorso image-based, Qwen2-VL-7B-Instruct MUST produrre una descrizione clinica testuale della lesione che rappresenti almeno forma, bordi, colore, texture e caratteristiche sospette richieste dal formato documentato.
\end{normbox}

\paragraph{SPO references (L2 support).}
\begin{lstlisting}
Qwen2-VL-7B-Instruct -- analyzes --> SkinLesionImage
Qwen2-VL-7B-Instruct -- produces --> ClinicalDescription
\end{lstlisting}

\subsubsection{FR-18 - Retrieval di casi storici con ChromaDB e all-MiniLM-L6-v2}\label{fr-18}
\begin{frmeta}
\begin{tabularx}{\linewidth}{L{43mm}Y}
\toprule
\textbf{ID} & \texttt{FR-18} \\
\textbf{Lifecycle} & \texttt{current} \\
\textbf{Authority} & \candidateauthority \\
\textbf{Type} & \FR \\
\textbf{Title} & Retrieval di casi storici con ChromaDB e all-MiniLM-L6-v2 \\
\textbf{Parent Decision} & \texttt{DEC-12} \\
\textbf{Owning MR (derived)} & \texttt{MR-01} \\
\textbf{SpecializedRequirement relation} & 0 active instances \\
\bottomrule
\end{tabularx}
\end{frmeta}

\paragraph{Functional obligation.}
Recuperare casi dermatologici storici simili alla descrizione clinica del caso corrente e renderli disponibili alla sintesi successiva.

\begin{normbox}[title={Functional clauses / inherited \texttt{normativeClause [1..*]}}]
DermaTriage MUST indicizzare e interrogare in ChromaDB le descrizioni dei casi usando embedding \nolinkurl{sentence-transformers/all-MiniLM-L6-v2}; per il caso corrente MUST calcolare la similarità cosine e recuperare i cinque casi storici più simili da usare come contesto della sintesi AI.
\end{normbox}

\paragraph{SPO references (L2 support).}
\begin{lstlisting}
ClinicalDescription -- embedded by --> all-MiniLM-L6-v2
ChromaDB -- retrieves by cosine similarity --> HistoricalCaseContext
HistoricalCaseContext -- consumed by --> BioMistral-7B
\end{lstlisting}

\subsubsection{FR-19 - Sintesi multi-source del triage con BioMistral-7B}\label{fr-19}
\begin{frmeta}
\begin{tabularx}{\linewidth}{L{43mm}Y}
\toprule
\textbf{ID} & \texttt{FR-19} \\
\textbf{Lifecycle} & \texttt{current} \\
\textbf{Authority} & \candidateauthority \\
\textbf{Type} & \FR \\
\textbf{Title} & Sintesi multi-source del triage con BioMistral-7B \\
\textbf{Parent Decision} & \texttt{DEC-12} \\
\textbf{Owning MR (derived)} & \texttt{MR-01} \\
\textbf{SpecializedRequirement relation} & 0 active instances \\
\bottomrule
\end{tabularx}
\end{frmeta}

\paragraph{Functional obligation.}
Combinare le informazioni prodotte dai precedenti stadi con i sintomi del caso e produrre la sintesi AI di triage.

\begin{normbox}[title={Functional clauses / inherited \texttt{normativeClause [1..*]}}]
BioMistral-7B MUST combinare, per lo stesso caso, la classificazione image-based, la descrizione clinica, il contesto dei casi storici recuperati e le informazioni sintomatologiche disponibili, producendo una sintesi contenente almeno urgency, confidence, reasoning e recommended action. Il campo \texttt{predicted\_pathology} è prodotto dal sistema ma non acquisisce per questo autorità di diagnosi clinica definitiva.
\end{normbox}

\paragraph{SPO references (L2 support).}
\begin{lstlisting}
BioMistral-7B -- consumes --> ImageUrgencyClassification
BioMistral-7B -- consumes --> ClinicalDescription
BioMistral-7B -- consumes --> HistoricalCaseContext
BioMistral-7B -- consumes --> AvailableSymptomInformation
BioMistral-7B -- produces --> AITriageSynthesis
\end{lstlisting}

\subsection{DEC-13 - Strategia di training della baseline EfficientNet-B4}\label{dec-13}
\begin{decmeta}
\begin{tabularx}{\linewidth}{L{43mm}Y}
\toprule
\textbf{ID} & \texttt{DEC-13} \\
\textbf{Lifecycle} & \texttt{current} \\
\textbf{Authority} & \candidateauthority \\
\textbf{Type} & \DEC \\
\textbf{Title} & Strategia di training della baseline EfficientNet-B4 \\
\textbf{Parent MacroRequirement} & \texttt{MR-01} \\
\textbf{Decision kind} & \texttt{SELECTABLE} \\
\bottomrule
\end{tabularx}
\end{decmeta}

\subsubsection{Context}
La capacità image-based di \texttt{FR-16} dipende da una baseline classificatoria addestrata e selezionata prima dell'uso operativo.

\subsubsection{Decision}
La baseline usa EfficientNet-B4 inizializzato da pesi ImageNet e viene addestrata su un training set bilanciato con SMOTE; validation e test rimangono split stratificati. Il training usa Adam, LinearWarmup seguito da CosineAnnealing, early stopping e selezione del best checkpoint tramite validation Macro F1.

\subsubsection{Consequences}
Il processo di costruzione della baseline resta distinto dal feedback-driven retraining del classificatore basato sulle successive correzioni cliniche.

\subsubsection{FR-20 - Addestramento della baseline classificatoria}\label{fr-20}
\begin{frmeta}
\begin{tabularx}{\linewidth}{L{43mm}Y}
\toprule
\textbf{ID} & \texttt{FR-20} \\
\textbf{Lifecycle} & \texttt{current} \\
\textbf{Authority} & \candidateauthority \\
\textbf{Type} & \FR \\
\textbf{Title} & Addestramento della baseline classificatoria \\
\textbf{Parent Decision} & \texttt{DEC-13} \\
\textbf{Owning MR (derived)} & \texttt{MR-01} \\
\textbf{SpecializedRequirement relation} & 0 active instances \\
\bottomrule
\end{tabularx}
\end{frmeta}

\paragraph{Functional obligation.}
Produrre la baseline EfficientNet-B4 secondo il processo di training e selezione del checkpoint stabilito dal progetto.

\begin{normbox}[title={Functional clauses / inherited \texttt{normativeClause [1..*]}}]
Il processo di training della baseline MUST usare un training set SMOTE-balanced da 2.286 righe, con 762 esempi per classe, e validation/test split stratificati da 345 righe; MUST usare Adam con learning rate \texttt{0.0001}, weight decay \texttt{1e-4}, warmup di 5 epoch, CosineAnnealing, massimo 50 epoch, early stopping patience 10, batch size 16, label smoothing 0.1, CrossEntropyLoss e AMP, selezionando il best checkpoint secondo validation Macro F1.
\end{normbox}

\paragraph{SPO references (L2 support).}
\begin{lstlisting}
TrainingDataset -- balanced by --> SMOTE
TrainingProcess -- trains --> EfficientNet-B4
ValidationMacroF1 -- selects --> BestBaselineCheckpoint
\end{lstlisting}

\subsection{DEC-14 - Esposizione del servizio e integrazione B4}\label{dec-14}
\begin{decmeta}
\begin{tabularx}{\linewidth}{L{43mm}Y}
\toprule
\textbf{ID} & \texttt{DEC-14} \\
\textbf{Lifecycle} & \texttt{current} \\
\textbf{Authority} & \candidateauthority \\
\textbf{Type} & \DEC \\
\textbf{Title} & Esposizione del servizio e integrazione B4 \\
\textbf{Parent MacroRequirement} & \texttt{MR-01} \\
\textbf{Decision kind} & \texttt{SELECTABLE} \\
\bottomrule
\end{tabularx}
\end{decmeta}

\subsubsection{Context}
DermaTriage supporta sia l'invocazione diretta del servizio sia un workflow in cui il caso proviene dal sistema esterno B4. B4 partecipa allo scambio di dati ma non acquisisce la responsabilità del triage posseduta da DermaTriage.

\subsubsection{Decision}
Il servizio DermaTriage è esposto tramite FastAPI e Uvicorn. Il percorso diretto usa \texttt{POST /analyze}; il percorso integrato con B4 usa \texttt{POST /diagnose} e consuma consultation, documenti e immagini resi disponibili da B4 per produrre e restituire l'esito di triage.

\subsubsection{Consequences}
Il boundary B4/DermaTriage resta esplicito e distinguibile dal percorso diretto. L'ambiente documentato usa Python 3.12.3 e PyTorch 2.11.x, con esecuzione NVIDIA L40S/CUDA e CPU fallback.

\subsubsection{FR-21 - Analisi diretta tramite POST /analyze}\label{fr-21}
\begin{frmeta}
\begin{tabularx}{\linewidth}{L{43mm}Y}
\toprule
\textbf{ID} & \texttt{FR-21} \\
\textbf{Lifecycle} & \texttt{current} \\
\textbf{Authority} & \candidateauthority \\
\textbf{Type} & \FR \\
\textbf{Title} & Analisi diretta tramite POST /analyze \\
\textbf{Parent Decision} & \texttt{DEC-14} \\
\textbf{Owning MR (derived)} & \texttt{MR-01} \\
\textbf{SpecializedRequirement relation} & 0 active instances \\
\bottomrule
\end{tabularx}
\end{frmeta}

\paragraph{Functional obligation.}
Consentire l'analisi diretta di un caso dermatologico senza dipendere dal workflow B4.

\begin{normbox}[title={Functional clauses / inherited \texttt{normativeClause [1..*]}}]
DermaTriage MUST esporre tramite FastAPI il percorso \texttt{POST /analyze} per sottoporre direttamente un caso al processo di triage supportato dal servizio.
\end{normbox}

\paragraph{SPO references (L2 support).}
\begin{lstlisting}
DirectClient -- invokes --> POST /analyze
POST /analyze -- invokes --> DermaTriageTriageProcess
\end{lstlisting}

\subsubsection{FR-22 - Diagnosi integrata B4 tramite POST /diagnose}\label{fr-22}
\begin{frmeta}
\begin{tabularx}{\linewidth}{L{43mm}Y}
\toprule
\textbf{ID} & \texttt{FR-22} \\
\textbf{Lifecycle} & \texttt{current} \\
\textbf{Authority} & \candidateauthority \\
\textbf{Type} & \FR \\
\textbf{Title} & Diagnosi integrata B4 tramite POST /diagnose \\
\textbf{Parent Decision} & \texttt{DEC-14} \\
\textbf{Owning MR (derived)} & \texttt{MR-01} \\
\textbf{SpecializedRequirement relation} & 0 active instances \\
\bottomrule
\end{tabularx}
\end{frmeta}

\paragraph{Functional obligation.}
Eseguire il triage di una consultation B4 utilizzando i dati resi disponibili dal sistema esterno e mantenendo la correlazione con il caso B4.

\begin{normbox}[title={Functional clauses / inherited \texttt{normativeClause [1..*]}}]
Quando viene invocato \texttt{POST /diagnose} per una consultation B4, DermaTriage MUST acquisire da B4 le informazioni e i documenti applicabili al caso, eseguire il processo di triage e mantenere l'associazione dell'esito con la stessa consultation. B4 resta un servizio/sistema esterno consumato e non il subject proprietario della valutazione di triage.
\end{normbox}

\paragraph{SPO references (L2 support).}
\begin{lstlisting}
B4 -- provides --> ConsultationData
B4 -- provides --> CaseDocuments
DermaTriage -- consumes --> ConsultationData
DermaTriage -- produces --> B4CorrelatedTriageOutcome
\end{lstlisting}

\subsection{DEC-01 - Continuità del triage in assenza di immagine}\label{dec-01}
\begin{decmeta}
\begin{tabularx}{\linewidth}{L{43mm}Y}
\toprule
\textbf{ID} & \texttt{DEC-01} \\
\textbf{Lifecycle} & \texttt{current} \\
\textbf{Authority} & \candidateauthority \\
\textbf{Type} & \DEC \\
\textbf{Title} & Continuità del triage in assenza di immagine \\
\textbf{Parent MacroRequirement} & \texttt{MR-01} \\
\textbf{Decision kind} & \texttt{SELECTABLE} \\
\bottomrule
\end{tabularx}
\end{decmeta}

\subsubsection{Context}
Il caso dermatologico può essere disponibile con o senza immagine. Quando l'immagine non è disponibile, la documentazione prevede un processo di \textbf{symptom-only scoring}; nel workflow B4-integrated le informazioni sintomatologiche provengono dai \textbf{B4 chatbot interaction fields}, tra cui itching, bleeding, growth/change e pain.

\subsubsection{Decision}
DermaTriage consente alla valutazione di urgenza di proseguire in assenza di immagine, utilizzando le informazioni sintomatologiche disponibili per il caso.

\subsubsection{Consequences}
L'immagine non è una precondizione assoluta per il triage. Il processo symptom-only è distinto dalla pipeline image-based di \texttt{DEC-12}; la completezza del contratto dei campi sintomatologici resta un gap documentale.

\subsubsection{FR-01 - Determinazione dell'urgenza su base sintomatologica in assenza di immagine}\label{fr-01}
\begin{frmeta}
\begin{tabularx}{\linewidth}{L{43mm}Y}
\toprule
\textbf{ID} & \texttt{FR-01} \\
\textbf{Lifecycle} & \texttt{current} \\
\textbf{Authority} & \candidateauthority \\
\textbf{Type} & \FR \\
\textbf{Title} & Determinazione dell'urgenza su base sintomatologica in assenza di immagine \\
\textbf{Parent Decision} & \texttt{DEC-01} \\
\textbf{Owning MR (derived)} & \texttt{MR-01} \\
\textbf{SpecializedRequirement relation} & 0 active instances \\
\bottomrule
\end{tabularx}
\end{frmeta}

\paragraph{Functional obligation.}
Determinare l'urgenza del caso quando l'immagine non è disponibile usando le informazioni sintomatologiche disponibili.

\begin{normbox}[title={Functional clauses / inherited \texttt{normativeClause [1..*]}}]
Quando per il caso dermatologico non è disponibile un'immagine, DermaTriage MUST determinare l'urgenza sulla base delle informazioni sintomatologiche disponibili. Nel workflow B4-integrated tali informazioni MUST essere derivate dai campi sintomatologici associati alla stessa consultation B4.
\end{normbox}

\paragraph{SPO references (L2 support).}
\begin{lstlisting}
AvailableSymptomInformation -- concerns --> DermatologicalCase
B4ChatbotInteractionFields -- provide --> AvailableSymptomInformation
DermaTriage -- determines --> SymptomBasedUrgency
\end{lstlisting}

\subsection{DEC-02 - Adozione della P-scale per la priorità operativa}\label{dec-02}
\begin{decmeta}
\begin{tabularx}{\linewidth}{L{43mm}Y}
\toprule
\textbf{ID} & \texttt{DEC-02} \\
\textbf{Lifecycle} & \texttt{current} \\
\textbf{Authority} & \candidateauthority \\
\textbf{Type} & \DEC \\
\textbf{Title} & Adozione della P-scale per la priorità operativa \\
\textbf{Parent MacroRequirement} & \texttt{MR-01} \\
\textbf{Decision kind} & \texttt{SELECTABLE} \\
\bottomrule
\end{tabularx}
\end{decmeta}

\subsubsection{Context}
La valutazione di urgenza deve essere trasformata in una priorità operativa leggibile dal processo di presa in carico. DermaTriage usa la scala P1-P4 e distingue il caso HIGH ad alta confidence dagli altri casi HIGH.

\subsubsection{Decision}
DermaTriage rappresenta la priorità operativa di triage mediante la P-scale P1-P4 e adotta il mapping HIGH con confidence maggiore di 0.85 verso P1, HIGH negli altri casi verso P2, MEDIUM verso P3 e LOW verso P4.

\subsubsection{Consequences}
Il comparatore per P1 è strettamente \texttt{> 0.85}. La documentazione associa anche P1=24h, P2=48h, P3=72h e P4=7 giorni, ma trigger, owner e natura normativa degli SLA restano non sufficientemente specificati. Il campo specialist è prodotto nel flusso operativo ma la regola completa di specialist selection appartiene a \texttt{MR-02} e resta aperta.

\subsubsection{FR-02 - Derivazione della priorità P-scale dalla valutazione di urgenza}\label{fr-02}
\begin{frmeta}
\begin{tabularx}{\linewidth}{L{43mm}Y}
\toprule
\textbf{ID} & \texttt{FR-02} \\
\textbf{Lifecycle} & \texttt{current} \\
\textbf{Authority} & \candidateauthority \\
\textbf{Type} & \FR \\
\textbf{Title} & Derivazione della priorità P-scale dalla valutazione di urgenza \\
\textbf{Parent Decision} & \texttt{DEC-02} \\
\textbf{Owning MR (derived)} & \texttt{MR-01} \\
\textbf{SpecializedRequirement relation} & 0 active instances \\
\bottomrule
\end{tabularx}
\end{frmeta}

\paragraph{Functional obligation.}
Derivare una priorità operativa P1-P4 dalla valutazione di urgenza e dalla confidence applicabile al caso.

\begin{normbox}[title={Functional clauses / inherited \texttt{normativeClause [1..*]}}]
\begin{tabularx}{\linewidth}{L{55mm}Y}
\toprule
\textbf{Condizione governata} & \textbf{Priorità richiesta} \\
\midrule
\texttt{urgency = HIGH} e \texttt{confidence > 0.85} & DermaTriage MUST produrre \texttt{P1}. \\
\texttt{urgency = HIGH} negli altri casi applicabili & DermaTriage MUST produrre \texttt{P2}. \\
\texttt{urgency = MEDIUM} & DermaTriage MUST produrre \texttt{P3}. \\
\texttt{urgency = LOW} & DermaTriage MUST produrre \texttt{P4}. \\
\bottomrule
\end{tabularx}
\end{normbox}

\paragraph{SPO references (L2 support).}
\begin{lstlisting}
UrgencyEvaluation -- contributes to --> OperationalTriagePriorityPScale
ApplicableConfidence -- refines HIGH mapping to --> P1
DermaTriage -- derives --> OperationalTriagePriorityPScale
\end{lstlisting}

\section{MR-02 - Indirizzamento specialistico}\label{mr-02}

\begin{mrmeta}
\begin{tabularx}{\linewidth}{L{43mm}Y}
\toprule
\textbf{ID} & \texttt{MR-02} \\
\textbf{Lifecycle} & \texttt{current} \\
\textbf{Authority} & \candidateauthority \\
\textbf{Type} & \MR \\
\textbf{Title} & Indirizzamento specialistico \\
\bottomrule
\end{tabularx}
\end{mrmeta}

\subsection{Intent}
Determinare un'indicazione di destinazione specialistica per il caso dermatologico, al fine di supportarne il successivo instradamento verso l'assistenza appropriata.

\subsection{Context}
La documentazione sorgente indica che DermaTriage produce o supporta un'informazione orientata alla destinazione specialistica. Nell'output operativo corrente il mapping P-scale include un campo \texttt{specialist} insieme allo SLA. Il significato disponibile non stabilisce però in modo sufficiente la regola di selezione, il vocabolario completo delle destinazioni né la responsabilità di assegnazione o prenotazione.

\subsection{Stakeholders}
Ruoli materialmente coinvolti:
\begin{itemize}
\item professionisti sanitari che utilizzano l'indicazione di destinazione;
\item servizi o unità specialistiche destinatarie del successivo instradamento;
\item soggetti cui il caso dermatologico si riferisce.
\end{itemize}

\subsection{Scope}
\textbf{In scope:} responsabilità macro di produrre/supportare un'indicazione di destinazione specialistica.

\textbf{Out of scope:} prenotazione, assegnazione effettiva dello specialista, completa tassonomia delle specialità, regola di fallback e relazione completa con la P-scale finché tali significati non sono governati.

\subsection{Assumptions / Constraints}
Nessuna regola di selezione specialistica completa viene assunta o ricostruita per inferenza. L'assenza di ulteriore decomposizione è intenzionale nel baseline corrente.

\subsection{dependsOn}
Nessuna dipendenza macro canonica è stabilita per \texttt{MR-02} nel baseline corrente.

\begin{statusbox}[title={STOP AT MR}]
Nella baseline corrente non vengono create Decision o FunctionalRequirement sotto \texttt{MR-02}. L'arresto preserva il significato disponibile senza riempire artificialmente la gerarchia.
\end{statusbox}


\section{MR-03 - Gestione della validazione clinica degli esiti}\label{mr-03}

\begin{mrmeta}
\begin{tabularx}{\linewidth}{L{43mm}Y}
\toprule
\textbf{ID} & \texttt{MR-03} \\
\textbf{Lifecycle} & \texttt{current} \\
\textbf{Authority} & \candidateauthority \\
\textbf{Type} & \MR \\
\textbf{Title} & Gestione della validazione clinica degli esiti \\
\bottomrule
\end{tabularx}
\end{mrmeta}

\subsection{Intent}
Supportare la registrazione e l'utilizzo della validazione o correzione effettuata da un professionista sanitario sugli esiti prodotti da DermaTriage, mantenendo distinguibile il risultato della revisione dall'esito originario.

\subsection{Context}
Il giudizio clinico appartiene al professionista sanitario. DermaTriage possiede la responsabilità di gestire, registrare e correlare l'esito della revisione al risultato originario del sistema, senza sostituirsi all'autorità clinica. Il progetto supporta sia operazioni amministrative DermaTriage sia scambio di validazioni con il sistema esterno B4.

\subsection{Stakeholders}
Ruoli materialmente coinvolti:
\begin{itemize}
\item professionisti sanitari che effettuano la revisione;
\item responsabili del servizio DermaTriage che utilizzano il risultato della revisione;
\item processi downstream di adattamento che consumano evidence derivata dalla revisione.
\end{itemize}

\subsection{Scope}
\textbf{In scope:} registrazione correlata del risultato della revisione; distinzione tra conferma e correzione; mantenimento della distinguibilità tra esito originario e revisione.

\textbf{Out of scope:} decisione clinica stessa; contenuto completo del payload di correzione; policy completa di history/retention/finality; semantica completa del disaccordo usato downstream.

\subsection{Assumptions / Constraints}
La revisione clinica è trattata come input proveniente da un professionista sanitario; il progetto non assume che conferma, correzione e disaccordo siano sinonimi.

\subsection{dependsOn}
Nessuna dipendenza macro canonica è stabilita per \texttt{MR-03} nel baseline corrente.

\subsection{DEC-03 - Separazione tra esito originario e revisione clinica}\label{dec-03}

\begin{decmeta}
\begin{tabularx}{\linewidth}{L{43mm}Y}
\toprule
\textbf{ID} & \texttt{DEC-03} \\
\textbf{Lifecycle} & \texttt{current} \\
\textbf{Authority} & \candidateauthority \\
\textbf{Type} & \DEC \\
\textbf{Title} & Separazione tra esito originario e revisione clinica \\
\textbf{Parent MacroRequirement} & \texttt{MR-03} \\
\textbf{Decision kind} & \texttt{SELECTABLE} \\
\bottomrule
\end{tabularx}
\end{decmeta}

\subsubsection{Context}
Una revisione clinica può confermare o correggere un esito prodotto da DermaTriage. Sovrascrivere semanticamente l'esito originario renderebbe indistinguibili il comportamento del sistema e il successivo giudizio clinico.

\subsubsection{Decision}
DermaTriage mantiene l'esito originario distinguibile dal successivo risultato della revisione clinica.

\subsubsection{Consequences}
L'esito originario e il risultato della revisione clinica restano distinguibili. Una conferma dell'esito originario resta distinguibile da una sua correzione.

\subsubsection{FR-03 - Registrazione correlata della revisione clinica}\label{fr-03}

\begin{frmeta}
\begin{tabularx}{\linewidth}{L{43mm}Y}
\toprule
\textbf{ID} & \texttt{FR-03} \\
\textbf{Lifecycle} & \texttt{current} \\
\textbf{Authority} & \candidateauthority \\
\textbf{Type} & \FR \\
\textbf{Title} & Registrazione correlata della revisione clinica \\
\textbf{Parent Decision} & \texttt{DEC-03} \\
\textbf{Owning MR (derived)} & \texttt{MR-03} \\
\textbf{SpecializedRequirement relation} & 0 active instances \\
\bottomrule
\end{tabularx}
\end{frmeta}

\paragraph{Functional obligation.}
Registrare il risultato della revisione clinica mantenendone la correlazione con l'esito DermaTriage originario.

\begin{normbox}[title={Functional clauses / inherited \texttt{normativeClause [1..*]}}]
Quando viene fornita una validazione o correzione clinica relativa a un esito DermaTriage, DermaTriage MUST registrare il risultato della revisione associandolo all'esito originario cui la revisione si riferisce.
\end{normbox}

\paragraph{SPO references (L2 support).}
\begin{lstlisting}
DermaTriage -- record --> ClinicalReviewResult
ClinicalReviewResult -- associated with --> OriginalDermaTriageOutcome
\end{lstlisting}



\subsubsection{FR-12 - Distinzione tra conferma e correzione della revisione clinica}\label{fr-12}

\begin{frmeta}
\begin{tabularx}{\linewidth}{L{43mm}Y}
\toprule
\textbf{ID} & \texttt{FR-12} \\
\textbf{Lifecycle} & \texttt{current} \\
\textbf{Authority} & \candidateauthority \\
\textbf{Type} & \FR \\
\textbf{Title} & Distinzione tra conferma e correzione della revisione clinica \\
\textbf{Parent Decision} & \texttt{DEC-03} \\
\textbf{Owning MR (derived)} & \texttt{MR-03} \\
\textbf{SpecializedRequirement relation} & 0 active instances \\
\bottomrule
\end{tabularx}
\end{frmeta}

\paragraph{Functional obligation.}
Rendere distinguibile se il risultato della revisione clinica conferma oppure corregge l'esito originario.

\begin{normbox}[title={Functional clauses / inherited \texttt{normativeClause [1..*]}}]
Per ogni risultato di revisione clinica registrato, DermaTriage MUST mantenere distinguibile una conferma dell'esito originario da una correzione dell'esito originario.
\end{normbox}

\paragraph{SPO references (L2 support).}
\begin{lstlisting}
ClinicalReviewResult -- classify disposition as --> ConfirmationOrCorrection
DermaTriage -- preserve distinction of --> ClinicalReviewDisposition
\end{lstlisting}





\subsection{DEC-15 - Scambio della validazione clinica con B4}\label{dec-15}
\begin{decmeta}
\begin{tabularx}{\linewidth}{L{43mm}Y}
\toprule
\textbf{ID} & \texttt{DEC-15} \\
\textbf{Lifecycle} & \texttt{current} \\
\textbf{Authority} & \candidateauthority \\
\textbf{Type} & \DEC \\
\textbf{Title} & Scambio della validazione clinica con B4 \\
\textbf{Parent MacroRequirement} & \texttt{MR-03} \\
\textbf{Decision kind} & \texttt{SELECTABLE} \\
\bottomrule
\end{tabularx}
\end{decmeta}

\subsubsection{Context}
Nel workflow integrato, la review clinica deve rimanere correlata alla consultation e all'esito di triage scambiati con B4.

\subsubsection{Decision}
DermaTriage usa le API B4 per il diagnostic-output write-back, il medical-validation write-back e il retrieval dei validated outcomes, mantenendo l'identità della consultation/case durante lo scambio.

\subsubsection{Consequences}
B4 resta un sistema esterno consumato; la responsabilità di registrare e correlare la review resta di DermaTriage.

\subsubsection{FR-23 - Write-back e retrieval della validazione B4}\label{fr-23}
\begin{frmeta}
\begin{tabularx}{\linewidth}{L{43mm}Y}
\toprule
\textbf{ID} & \texttt{FR-23} \\
\textbf{Lifecycle} & \texttt{current} \\
\textbf{Authority} & \candidateauthority \\
\textbf{Type} & \FR \\
\textbf{Title} & Write-back e retrieval della validazione B4 \\
\textbf{Parent Decision} & \texttt{DEC-15} \\
\textbf{Owning MR (derived)} & \texttt{MR-03} \\
\textbf{SpecializedRequirement relation} & 0 active instances \\
\bottomrule
\end{tabularx}
\end{frmeta}

\paragraph{Functional obligation.}
Scambiare con B4 l'esito diagnostico/di triage e la successiva validazione mantenendo la correlazione con lo stesso caso.

\begin{normbox}[title={Functional clauses / inherited \texttt{normativeClause [1..*]}}]
Nel workflow B4-integrated, DermaTriage MUST mantenere la stessa consultation/case identity durante diagnostic-output write-back, medical-validation write-back e retrieval dei validated outcomes, così che la revisione resti associabile all'esito originario cui si riferisce.
\end{normbox}

\paragraph{SPO references (L2 support).}
\begin{lstlisting}
DermaTriage -- writes back --> B4TriageOutcome
ClinicianValidation -- written back to --> B4
B4 -- provides --> ValidatedOutcome
ValidatedOutcome -- correlatedWith --> OriginalDermaTriageOutcome
\end{lstlisting}

\subsection{DEC-16 - Autenticazione delle interfacce di review}\label{dec-16}
\begin{decmeta}
\begin{tabularx}{\linewidth}{L{43mm}Y}
\toprule
\textbf{ID} & \texttt{DEC-16} \\
\textbf{Lifecycle} & \texttt{current} \\
\textbf{Authority} & \candidateauthority \\
\textbf{Type} & \DEC \\
\textbf{Title} & Autenticazione delle interfacce di review \\
\textbf{Parent MacroRequirement} & \texttt{MR-03} \\
\textbf{Decision kind} & \texttt{SELECTABLE} \\
\bottomrule
\end{tabularx}
\end{decmeta}

\subsubsection{Context}
Le operazioni amministrative di review/correction e l'interazione con B4 attraversano boundary differenti e usano meccanismi di autenticazione distinti.

\subsubsection{Decision}
Le operazioni amministrative DermaTriage usano \texttt{X-API-Key}; l'interazione con B4 usa bearer JWT, con acquisizione e refresh del token gestiti dal client B4.

\subsubsection{Consequences}
Le operazioni amministrative DermaTriage e l'interazione con B4 usano meccanismi di autenticazione distinti.

\subsubsection{FR-24 - X-API-Key per operazioni amministrative DermaTriage}\label{fr-24}
\begin{frmeta}
\begin{tabularx}{\linewidth}{L{43mm}Y}
\toprule
\textbf{ID} & \texttt{FR-24} \\
\textbf{Lifecycle} & \texttt{current} \\
\textbf{Authority} & \candidateauthority \\
\textbf{Type} & \FR \\
\textbf{Title} & X-API-Key per operazioni amministrative DermaTriage \\
\textbf{Parent Decision} & \texttt{DEC-16} \\
\textbf{Owning MR (derived)} & \texttt{MR-03} \\
\textbf{SpecializedRequirement relation} & 0 active instances \\
\bottomrule
\end{tabularx}
\end{frmeta}
\paragraph{Functional obligation.} Richiedere la API key prevista per le operazioni amministrative protette.
\begin{normbox}[title={Functional clauses / inherited \texttt{normativeClause [1..*]}}]
Le operazioni amministrative DermaTriage di review, validation/correction, prompt management e classifier retraining esposte come protette MUST richiedere il meccanismo \texttt{X-API-Key} previsto dall'interfaccia.
\end{normbox}
\paragraph{SPO references (L2 support).}
\begin{lstlisting}
AdministrativeClient -- presents --> X-API-Key
ProtectedDermaTriageOperation -- requires --> X-API-Key
\end{lstlisting}

\subsubsection{FR-25 - Bearer JWT per interazione B4}\label{fr-25}
\begin{frmeta}
\begin{tabularx}{\linewidth}{L{43mm}Y}
\toprule
\textbf{ID} & \texttt{FR-25} \\
\textbf{Lifecycle} & \texttt{current} \\
\textbf{Authority} & \candidateauthority \\
\textbf{Type} & \FR \\
\textbf{Title} & Bearer JWT per interazione B4 \\
\textbf{Parent Decision} & \texttt{DEC-16} \\
\textbf{Owning MR (derived)} & \texttt{MR-03} \\
\textbf{SpecializedRequirement relation} & 0 active instances \\
\bottomrule
\end{tabularx}
\end{frmeta}
\paragraph{Functional obligation.} Autenticare le chiamate DermaTriage verso B4 mediante il flusso bearer JWT documentato.
\begin{normbox}[title={Functional clauses / inherited \texttt{normativeClause [1..*]}}]
Il client B4 di DermaTriage MUST acquisire e rinnovare il token previsto dal flusso di autenticazione B4 e MUST usare il bearer JWT nelle chiamate B4 che richiedono autenticazione.
\end{normbox}
\paragraph{SPO references (L2 support).}
\begin{lstlisting}
DermaTriageB4Client -- obtains/refreshes --> B4BearerJWT
DermaTriageB4Client -- presents --> B4BearerJWT
\end{lstlisting}

\section{MR-04 - Adattamento controllato sulla base della revisione clinica}\label{mr-04}

\begin{mrmeta}
\begin{tabularx}{\linewidth}{L{43mm}Y}
\toprule
\textbf{ID} & \texttt{MR-04} \\
\textbf{Lifecycle} & \texttt{current} \\
\textbf{Authority} & \candidateauthority \\
\textbf{Type} & \MR \\
\textbf{Title} & Adattamento controllato sulla base della revisione clinica \\
\bottomrule
\end{tabularx}
\end{mrmeta}

\subsection{Intent}
Consentire l'evoluzione controllata del comportamento di DermaTriage utilizzando evidence derivata dalle validazioni e correzioni cliniche accumulate, evitando che modifiche degradanti vengano accettate senza controllo.

\subsection{Context}
DermaTriage utilizza risultati di revisione clinica per supportare l'evoluzione di capability differenti, in particolare il percorso di evoluzione dei prompt usati da Qwen2-VL e BioMistral e il percorso di adattamento di EfficientNet-B4. I due percorsi hanno condizioni di attivazione, evidence qualificante e risultati di lifecycle distinti.

\subsection{Stakeholders}
Ruoli materialmente coinvolti:
\begin{itemize}
\item responsabili dell'evoluzione e manutenzione di DermaTriage;
\item professionisti sanitari le cui revisioni forniscono evidence clinica;
\item utilizzatori e soggetti influenzati dalla qualità del triage dopo l'adattamento.
\end{itemize}

\subsection{Scope}
\textbf{In scope:} attivazione su accumulo di evidence; separazione dei percorsi di adattamento; selezione dell'evidence recente; qualificazione del disaccordo clinico; derivazione del target di supervisione; valutazione comparativa; proprietà qualitative di accettazione; reversibilità post-adoption.

\textbf{Out of scope:} giudizio clinico posseduto da \texttt{MR-03}; autorità finale di deployment non governata.

\subsection{Assumptions / Constraints}
Le revisioni cliniche consumate da questo branch sono governate da \texttt{MR-03}. I percorsi di adattamento possono riusare dati solo quando tali dati soddisfano indipendentemente le regole di ciascun percorso.

\subsection{dependsOn}
\texttt{MR-04 dependsOn MR-03}. La dipendenza esprime consumo del significato di revisione clinica posseduto da \texttt{MR-03}; non crea co-ownership dei Requirement.

\subsection{DEC-04 - Attivazione dell'adattamento su accumulo di evidence clinica}\label{dec-04}

\begin{decmeta}
\begin{tabularx}{\linewidth}{L{43mm}Y}
\toprule
\textbf{ID} & \texttt{DEC-04} \\
\textbf{Lifecycle} & \texttt{current} \\
\textbf{Authority} & \candidateauthority \\
\textbf{Type} & \DEC \\
\textbf{Title} & Attivazione dell'adattamento su accumulo di evidence clinica \\
\textbf{Parent MacroRequirement} & \texttt{MR-04} \\
\textbf{Decision kind} & \texttt{SELECTABLE} \\
\bottomrule
\end{tabularx}
\end{decmeta}

\subsubsection{Context}
L'adattamento non viene attivato per ogni singola revisione clinica. I diversi percorsi devono accumulare evidence pertinente prima di avviare il proprio ciclo di evoluzione.

\subsubsection{Decision}
Ogni processo di adattamento viene attivato soltanto quando l'evidence clinica pertinente accumulata soddisfa la condizione di accumulo governata per quel processo.

\subsubsection{Consequences}
Il percorso prompt e il percorso classificatorio hanno soglie distinte: 10 correzioni pertinenti per il prompt e 50 correzioni qualificanti per il classifier. \texttt{FR-04} e \texttt{FR-05} governano separatamente i due trigger.

\subsubsection{FR-04 - Attivazione del percorso di evoluzione del prompt}\label{fr-04}

\begin{frmeta}
\begin{tabularx}{\linewidth}{L{43mm}Y}
\toprule
\textbf{ID} & \texttt{FR-04} \\
\textbf{Lifecycle} & \texttt{current} \\
\textbf{Authority} & \candidateauthority \\
\textbf{Type} & \FR \\
\textbf{Title} & Attivazione del percorso di evoluzione del prompt \\
\textbf{Parent Decision} & \texttt{DEC-04} \\
\textbf{Owning MR (derived)} & \texttt{MR-04} \\
\textbf{SpecializedRequirement relation} & 0 active instances \\
\bottomrule
\end{tabularx}
\end{frmeta}

\paragraph{Functional obligation.}
Attivare un ciclo di evoluzione del prompt quando l'evidence di correzione clinica pertinente raggiunge la soglia applicabile al percorso prompt.

\begin{normbox}[title={Functional clauses / inherited \texttt{normativeClause [1..*]}}]
Quando l'evidence di correzione clinica pertinente accumulata raggiunge 10 elementi, DermaTriage MUST attivare un ciclo di evoluzione del prompt.
\end{normbox}

\paragraph{SPO references (L2 support).}
\begin{lstlisting}
AccumulatedPromptCorrectionEvidence -- reaches --> PromptEvolutionThreshold
DermaTriage -- activate --> PromptEvolutionCycle
\end{lstlisting}


\subsubsection{FR-05 - Attivazione del percorso di adattamento classificatorio}\label{fr-05}

\begin{frmeta}
\begin{tabularx}{\linewidth}{L{43mm}Y}
\toprule
\textbf{ID} & \texttt{FR-05} \\
\textbf{Lifecycle} & \texttt{current} \\
\textbf{Authority} & \candidateauthority \\
\textbf{Type} & \FR \\
\textbf{Title} & Attivazione del percorso di adattamento classificatorio \\
\textbf{Parent Decision} & \texttt{DEC-04} \\
\textbf{Owning MR (derived)} & \texttt{MR-04} \\
\textbf{SpecializedRequirement relation} & 0 active instances \\
\bottomrule
\end{tabularx}
\end{frmeta}

\paragraph{Functional obligation.}
Attivare il ciclo di adattamento della classificazione quando l'evidence qualificante accumulata raggiunge la soglia applicabile al percorso classificatorio.

\begin{normbox}[title={Functional clauses / inherited \texttt{normativeClause [1..*]}}]
Quando l'evidence qualificante accumulata raggiunge 50 correzioni, DermaTriage MUST attivare il ciclo di adattamento classificatorio di EfficientNet-B4.
\end{normbox}

\paragraph{SPO references (L2 support).}
\begin{lstlisting}
AccumulatedClassifierEvidence -- reaches --> ClassifierAdaptationThreshold
DermaTriage -- activate --> ClassifierAdaptationCycle
\end{lstlisting}


\paragraph{Cross-branch / consumption note.}
\texttt{FR-05} consuma evidence di disaccordo qualificata secondo \texttt{FR-07}. Questo consumo non modifica la parent Decision di nessuno dei due Requirement.

\subsection{DEC-05 - Separazione dei percorsi di adattamento per capability}\label{dec-05}

\begin{decmeta}
\begin{tabularx}{\linewidth}{L{43mm}Y}
\toprule
\textbf{ID} & \texttt{DEC-05} \\
\textbf{Lifecycle} & \texttt{current} \\
\textbf{Authority} & \candidateauthority \\
\textbf{Type} & \DEC \\
\textbf{Title} & Separazione dei percorsi di adattamento per capability \\
\textbf{Parent MacroRequirement} & \texttt{MR-04} \\
\textbf{Decision kind} & \texttt{SELECTABLE} \\
\bottomrule
\end{tabularx}
\end{decmeta}

\subsubsection{Context}
L'evoluzione del prompt e l'adattamento della classificazione sono capability distinte. Condividere dati o trovarsi nello stesso progetto non rende equivalenti le loro condizioni di attivazione, le loro regole di qualificazione dell'evidence o i loro risultati di lifecycle.

\subsubsection{Decision}
DermaTriage governa i percorsi di adattamento come lifecycle distinti: condizioni di attivazione, evidence qualificante e risultati di qualificazione/accettazione non si trasferiscono automaticamente da un percorso all'altro.

\subsubsection{Consequences}
Le condizioni di attivazione, la qualificazione dell'evidence e i risultati di lifecycle restano indipendenti tra i diversi percorsi di adattamento.

\begin{statusbox}
\texttt{FR-11 - Indipendenza dei lifecycle di adattamento} è \textbf{superseded}. Il suo ID resta preservato come identità storica; la baseline corrente usa \texttt{FR-13}, \texttt{FR-14} e \texttt{FR-15} per obbligazioni indipendentemente revisionabili e assessable.
\end{statusbox}

\subsubsection{FR-13 - Indipendenza delle condizioni di attivazione tra percorsi di adattamento}\label{fr-13}

\begin{frmeta}
\begin{tabularx}{\linewidth}{L{43mm}Y}
\toprule
\textbf{ID} & \texttt{FR-13} \\
\textbf{Lifecycle} & \texttt{current} \\
\textbf{Authority} & \candidateauthority \\
\textbf{Type} & \FR \\
\textbf{Title} & Indipendenza delle condizioni di attivazione tra percorsi di adattamento \\
\textbf{Parent Decision} & \texttt{DEC-05} \\
\textbf{Owning MR (derived)} & \texttt{MR-04} \\
\textbf{SpecializedRequirement relation} & 0 active instances \\
\bottomrule
\end{tabularx}
\end{frmeta}

\paragraph{Functional obligation.}
Impedire che la condizione di attivazione soddisfatta da un percorso sia considerata, da sola, sufficiente ad attivare un altro percorso.

\begin{normbox}[title={Functional clauses / inherited \texttt{normativeClause [1..*]}}]
Il soddisfacimento della condizione di attivazione di un percorso di adattamento MUST NOT, di per sé, essere trattato come soddisfacimento della condizione di attivazione di un altro percorso di adattamento.
\end{normbox}

\paragraph{SPO references (L2 support).}
\begin{lstlisting}
ActivationConditionOfPathA -- MUST NOT imply --> ActivationConditionOfPathB
\end{lstlisting}

\subsubsection{FR-14 - Indipendenza della qualificazione dell'evidence tra percorsi di adattamento}\label{fr-14}

\begin{frmeta}
\begin{tabularx}{\linewidth}{L{43mm}Y}
\toprule
\textbf{ID} & \texttt{FR-14} \\
\textbf{Lifecycle} & \texttt{current} \\
\textbf{Authority} & \candidateauthority \\
\textbf{Type} & \FR \\
\textbf{Title} & Indipendenza della qualificazione dell'evidence tra percorsi di adattamento \\
\textbf{Parent Decision} & \texttt{DEC-05} \\
\textbf{Owning MR (derived)} & \texttt{MR-04} \\
\textbf{SpecializedRequirement relation} & 0 active instances \\
\bottomrule
\end{tabularx}
\end{frmeta}

\paragraph{Functional obligation.}
Impedire che evidence qualificante per un percorso sia considerata automaticamente qualificante per un altro percorso.

\begin{normbox}[title={Functional clauses / inherited \texttt{normativeClause [1..*]}}]
Evidence che qualifica per un percorso di adattamento MUST NOT, di per sé, essere trattata come evidence qualificante per un altro percorso di adattamento.
\end{normbox}

\paragraph{SPO references (L2 support).}
\begin{lstlisting}
EvidenceQualifiedForPathA -- MUST NOT imply --> EvidenceQualifiedForPathB
\end{lstlisting}

\subsubsection{FR-15 - Indipendenza dei risultati di qualificazione e accettazione tra percorsi di adattamento}\label{fr-15}

\begin{frmeta}
\begin{tabularx}{\linewidth}{L{43mm}Y}
\toprule
\textbf{ID} & \texttt{FR-15} \\
\textbf{Lifecycle} & \texttt{current} \\
\textbf{Authority} & \candidateauthority \\
\textbf{Type} & \FR \\
\textbf{Title} & Indipendenza dei risultati di qualificazione e accettazione tra percorsi di adattamento \\
\textbf{Parent Decision} & \texttt{DEC-05} \\
\textbf{Owning MR (derived)} & \texttt{MR-04} \\
\textbf{SpecializedRequirement relation} & 0 active instances \\
\bottomrule
\end{tabularx}
\end{frmeta}

\paragraph{Functional obligation.}
Impedire che un risultato di qualificazione, accettazione o progressione raggiunto da un percorso venga trattato automaticamente come risultato equivalente per un altro percorso.

\begin{normbox}[title={Functional clauses / inherited \texttt{normativeClause [1..*]}}]
Un risultato di qualificazione, accettazione o progressione raggiunto da un percorso di adattamento MUST NOT, di per sé, essere trattato come risultato di lifecycle equivalente per un altro percorso di adattamento.
\end{normbox}

\paragraph{SPO references (L2 support).}
\begin{lstlisting}
LifecycleResultOfPathA -- MUST NOT imply --> EquivalentLifecycleResultOfPathB
\end{lstlisting}

\subsection{DEC-06 - Accettazione comparativa dell'adattamento della classificazione di urgenza}\label{dec-06}

\begin{decmeta}
\begin{tabularx}{\linewidth}{L{43mm}Y}
\toprule
\textbf{ID} & \texttt{DEC-06} \\
\textbf{Lifecycle} & \texttt{current} \\
\textbf{Authority} & \candidateauthority \\
\textbf{Type} & \DEC \\
\textbf{Title} & Accettazione comparativa dell'adattamento della classificazione di urgenza \\
\textbf{Parent MacroRequirement} & \texttt{MR-04} \\
\textbf{Decision kind} & \texttt{SELECTABLE} \\
\bottomrule
\end{tabularx}
\end{decmeta}

\subsubsection{Context}
Un candidato di adattamento classificatorio non deve essere adottato soltanto perché è stato prodotto. La reference EfficientNet-B4 model version \texttt{1.0.0} riporta overall accuracy 81.74\%, Macro F1 0.8036, AUC-ROC 0.906, HIGH sensitivity 90.24\% e false LOW rate 7.93\%; la fonte riporta anche HIGH recall 81.71\% e 134/164 casi HIGH corretti, con relazione fra i due valori non chiarita. Serve una valutazione comparativa rispetto alla reference applicabile.

\subsubsection{Decision}
Un candidato di adattamento della classificazione di urgenza deve essere valutato comparativamente rispetto alla reference applicabile e può essere considerato qualificato per l'adozione soltanto quando soddisfa tutti i criteri comparativi governati applicabili.

\subsubsection{Consequences}
\texttt{FR-09} governa valutazione e qualificazione comparativa. La qualificazione non equivale a deployment automatico. \texttt{DEC-07} governa proprietà qualitative aggiuntive dei criteri di accettazione.

\subsubsection{FR-09 - Qualificazione comparativa del candidato di adattamento classificatorio}\label{fr-09}

\begin{frmeta}
\begin{tabularx}{\linewidth}{L{43mm}Y}
\toprule
\textbf{ID} & \texttt{FR-09} \\
\textbf{Lifecycle} & \texttt{current} \\
\textbf{Authority} & \candidateauthority \\
\textbf{Type} & \FR \\
\textbf{Title} & Qualificazione comparativa del candidato di adattamento classificatorio \\
\textbf{Parent Decision} & \texttt{DEC-06} \\
\textbf{Owning MR (derived)} & \texttt{MR-04} \\
\textbf{SpecializedRequirement relation} & 0 active instances in the current baseline. \texttt{DEC-07} preserves the applicable non-security quality meaning without introducing a concrete subtype. \\
\bottomrule
\end{tabularx}
\end{frmeta}

\paragraph{Functional obligation.}
Valutare comparativamente un candidato classificatorio e considerarlo qualificato per l'adozione solo se soddisfa tutti i criteri applicabili.

\begin{normbox}[title={Functional clauses / inherited \texttt{normativeClause [1..*]}}]
\textbf{NC-09.1.} Quando un candidato di adattamento della classificazione di urgenza viene considerato per l'adozione, DermaTriage MUST valutarlo comparativamente rispetto alla versione di riferimento applicabile usando i criteri di accettazione governati.

\medskip
\textbf{NC-09.2.} DermaTriage MUST considerare il candidato qualificato per l'adozione soltanto se tutti i criteri comparativi governati applicabili risultano soddisfatti.
\end{normbox}

\paragraph{SPO references (L2 support).}
\begin{lstlisting}
DermaTriage -- comparatively evaluate --> CandidateUrgencyClassificationAdaptation
CandidateUrgencyClassificationAdaptation -- compared against --> ApplicableReferenceVersion
SatisfiedAcceptanceCriteria -- qualify --> CandidateUrgencyClassificationAdaptation
\end{lstlisting}


\subsection{DEC-07 - Prioritizzazione asimmetrica delle metriche di qualità}\label{dec-07}

\begin{decmeta}
\begin{tabularx}{\linewidth}{L{43mm}Y}
\toprule
\textbf{ID} & \texttt{DEC-07} \\
\textbf{Lifecycle} & \texttt{current} \\
\textbf{Authority} & \candidateauthority \\
\textbf{Type} & \DEC \\
\textbf{Title} & Prioritizzazione asimmetrica delle metriche di qualità \\
\textbf{Parent MacroRequirement} & \texttt{MR-04} \\
\textbf{Decision kind} & \texttt{SELECTABLE} \\
\bottomrule
\end{tabularx}
\end{decmeta}

\subsubsection{Context}
La policy di accettazione non tratta tutte le metriche come intercambiabili. La documentazione attribuisce priorità alla non-degradazione di sensitivity e false-low, consentendo per l'accuracy soltanto una degradazione limitata.

\subsubsection{Decision}
Nel confronto con la reference applicabile, sensitivity e false-low performance non devono peggiorare, mentre l'overall accuracy può degradare al massimo del 5\%.

\subsubsection{Consequences}
La qualificazione comparativa tratta in modo asimmetrico le metriche: sensitivity e false-low performance non devono peggiorare, mentre l'overall accuracy può degradare soltanto entro la tolleranza applicabile.


\subsection{DEC-08 - Reversibilità dell'adattamento della classificazione in caso di degrado}\label{dec-08}

\begin{decmeta}
\begin{tabularx}{\linewidth}{L{43mm}Y}
\toprule
\textbf{ID} & \texttt{DEC-08} \\
\textbf{Lifecycle} & \texttt{current} \\
\textbf{Authority} & \candidateauthority \\
\textbf{Type} & \DEC \\
\textbf{Title} & Reversibilità dell'adattamento della classificazione in caso di degrado \\
\textbf{Parent MacroRequirement} & \texttt{MR-04} \\
\textbf{Decision kind} & \texttt{SELECTABLE} \\
\bottomrule
\end{tabularx}
\end{decmeta}

\subsubsection{Context}
Un adattamento già adottato può manifestare degrado materiale dopo l'adozione. La capacità di recuperare una versione o stato precedente deve restare distinta dalla qualificazione pre-adoption.

\subsubsection{Decision}
Un adattamento della classificazione di urgenza già adottato che manifesta un degrado di accuracy superiore al 5\% rispetto alla reference post-adoption applicabile deve poter essere revocato ripristinando una versione o stato precedente accettabile.

\subsubsection{Consequences}
\texttt{FR-10} governa la revoca/ripristino. Il requisito non stabilisce che il rollback sia automatico. Il threshold post-adoption \texttt{R} resta semanticamente distinto dalla tolleranza pre-adoption \texttt{T} anche quando i literal correnti coincidono.

\subsubsection{FR-10 - Revoca di un adattamento classificatorio degradante}\label{fr-10}

\begin{frmeta}
\begin{tabularx}{\linewidth}{L{43mm}Y}
\toprule
\textbf{ID} & \texttt{FR-10} \\
\textbf{Lifecycle} & \texttt{current} \\
\textbf{Authority} & \candidateauthority \\
\textbf{Type} & \FR \\
\textbf{Title} & Revoca di un adattamento classificatorio degradante \\
\textbf{Parent Decision} & \texttt{DEC-08} \\
\textbf{Owning MR (derived)} & \texttt{MR-04} \\
\textbf{SpecializedRequirement relation} & 0 active instances \\
\bottomrule
\end{tabularx}
\end{frmeta}

\paragraph{Functional obligation.}
Supportare la revoca di un adattamento classificatorio già adottato quando manifesta degrado materiale, ripristinando una versione o stato precedente accettabile.

\begin{normbox}[title={Functional clauses / inherited \texttt{normativeClause [1..*]}}]
Per un adattamento della classificazione di urgenza già adottato che manifesti un degrado di accuracy superiore al 5\% rispetto alla reference post-adoption applicabile, DermaTriage MUST supportarne la revoca mediante il ripristino di una versione o stato precedente accettabile; le versioni ripristinabili sono mantenute in \nolinkurl{models/versions/} con tracking in \nolinkurl{db/model_versions.json} e modello attivo in \nolinkurl{models/efficientnet_b4.pth}.
\end{normbox}

\paragraph{SPO references (L2 support).}
\begin{lstlisting}
MateriallyDegradingAdoptedAdaptation -- triggers support for --> Revocation
DermaTriage -- restore --> PreviousAcceptableVersionOrState
\end{lstlisting}


\subsection{DEC-09 - Uso di evidence recente e limitata nel percorso di evoluzione del prompt}\label{dec-09}

\begin{decmeta}
\begin{tabularx}{\linewidth}{L{43mm}Y}
\toprule
\textbf{ID} & \texttt{DEC-09} \\
\textbf{Lifecycle} & \texttt{current} \\
\textbf{Authority} & \candidateauthority \\
\textbf{Type} & \DEC \\
\textbf{Title} & Uso di evidence recente e limitata nel percorso di evoluzione del prompt \\
\textbf{Parent MacroRequirement} & \texttt{MR-04} \\
\textbf{Decision kind} & \texttt{SELECTABLE} \\
\bottomrule
\end{tabularx}
\end{decmeta}

\subsubsection{Context}
L'evoluzione del prompt non utilizza una storia illimitata delle correzioni cliniche. La fonte governa l'uso di evidence recente e limitata per costruire il set del ciclo corrente.

\subsubsection{Decision}
Per l'evoluzione del prompt, DermaTriage usa una finestra scorrevole e limitata delle correzioni cliniche pertinenti più recenti.

\subsubsection{Consequences}
\texttt{FR-06} governa la selezione dell'evidence set. La finestra contiene i 20 esempi pertinenti più recenti; ordering, deduplication, underfill e riuso tra cicli restano non sufficientemente specificati.

\subsubsection{FR-06 - Selezione dell'evidence recente per l'evoluzione del prompt}\label{fr-06}

\begin{frmeta}
\begin{tabularx}{\linewidth}{L{43mm}Y}
\toprule
\textbf{ID} & \texttt{FR-06} \\
\textbf{Lifecycle} & \texttt{current} \\
\textbf{Authority} & \candidateauthority \\
\textbf{Type} & \FR \\
\textbf{Title} & Selezione dell'evidence recente per l'evoluzione del prompt \\
\textbf{Parent Decision} & \texttt{DEC-09} \\
\textbf{Owning MR (derived)} & \texttt{MR-04} \\
\textbf{SpecializedRequirement relation} & 0 active instances \\
\bottomrule
\end{tabularx}
\end{frmeta}

\paragraph{Functional obligation.}
Costruire l'evidence set del ciclo prompt usando una finestra scorrevole limitata delle correzioni cliniche pertinenti più recenti.

\begin{normbox}[title={Functional clauses / inherited \texttt{normativeClause [1..*]}}]
Per ciascun ciclo di evoluzione del prompt, DermaTriage MUST costruire l'evidence set utilizzando una finestra scorrevole dei 20 esempi di correzione clinica pertinenti più recenti.
\end{normbox}

\paragraph{SPO references (L2 support).}
\begin{lstlisting}
DermaTriage -- select recent bounded --> PromptEvolutionEvidenceSet
RecentPertinentClinicalCorrections -- populate --> PromptEvolutionEvidenceSet
\end{lstlisting}


\subsection{DEC-10 - Qualificazione del disaccordo clinico come evidence per l'adattamento classificatorio}\label{dec-10}

\begin{decmeta}
\begin{tabularx}{\linewidth}{L{43mm}Y}
\toprule
\textbf{ID} & \texttt{DEC-10} \\
\textbf{Lifecycle} & \texttt{current} \\
\textbf{Authority} & \candidateauthority \\
\textbf{Type} & \DEC \\
\textbf{Title} & Qualificazione del disaccordo clinico come evidence per l'adattamento classificatorio \\
\textbf{Parent MacroRequirement} & \texttt{MR-04} \\
\textbf{Decision kind} & \texttt{SELECTABLE} \\
\bottomrule
\end{tabularx}
\end{decmeta}

\subsubsection{Context}
Non ogni revisione clinica è trattata come evidence equivalente per l'adattamento della classificazione di urgenza. La fonte distingue le revisioni che esprimono disaccordo rispetto all'esito AI pertinente.

\subsubsection{Decision}
Per il percorso di adattamento della classificazione di urgenza, DermaTriage tratta come evidence qualificante sulla dimensione di disaccordo soltanto le revisioni cliniche che esprimono disaccordo rispetto all'esito AI pertinente.

\subsubsection{Consequences}
\texttt{FR-07} governa la qualificazione. Il concetto governato è \texttt{ClinicianDisagreement}; l'encoding sorgente corrente \texttt{agrees == False} resta una rappresentazione di stato/dato e non l'identità normativa del concetto.

\subsubsection{FR-07 - Qualificazione del disaccordo clinico come evidence per l'adattamento classificatorio}\label{fr-07}

\begin{frmeta}
\begin{tabularx}{\linewidth}{L{43mm}Y}
\toprule
\textbf{ID} & \texttt{FR-07} \\
\textbf{Lifecycle} & \texttt{current} \\
\textbf{Authority} & \candidateauthority \\
\textbf{Type} & \FR \\
\textbf{Title} & Qualificazione del disaccordo clinico come evidence per l'adattamento classificatorio \\
\textbf{Parent Decision} & \texttt{DEC-10} \\
\textbf{Owning MR (derived)} & \texttt{MR-04} \\
\textbf{SpecializedRequirement relation} & 0 active instances \\
\bottomrule
\end{tabularx}
\end{frmeta}

\paragraph{Functional obligation.}
Qualificare come evidence per il percorso classificatorio soltanto le revisioni cliniche che esprimono disaccordo rispetto all'esito AI pertinente.

\begin{normbox}[title={Functional clauses / inherited \texttt{normativeClause [1..*]}}]
Nel determinare l'evidence clinica qualificante per il percorso di adattamento della classificazione di urgenza, DermaTriage MUST considerare qualificanti soltanto le revisioni cliniche che esprimono disaccordo rispetto all'esito AI pertinente; nel payload B4 documentato tale stato è rappresentato da \texttt{agrees == False}.
\end{normbox}

\paragraph{SPO references (L2 support).}
\begin{lstlisting}
ClinicalReview -- expresses --> ClinicianDisagreement
ClinicianDisagreement -- qualifies as --> ClassifierAdaptationEvidence
\end{lstlisting}


\paragraph{Cross-branch / consumption note.}
L'evidence qualificata da \texttt{FR-07} viene consumata da \texttt{FR-05} per l'accumulo che conduce all'attivazione del percorso classificatorio.

\subsection{DEC-11 - Derivazione della supervision classificatoria dalla priorità clinicamente corretta}\label{dec-11}

\begin{decmeta}
\begin{tabularx}{\linewidth}{L{43mm}Y}
\toprule
\textbf{ID} & \texttt{DEC-11} \\
\textbf{Lifecycle} & \texttt{current} \\
\textbf{Authority} & \candidateauthority \\
\textbf{Type} & \DEC \\
\textbf{Title} & Derivazione della supervision classificatoria dalla priorità clinicamente corretta \\
\textbf{Parent MacroRequirement} & \texttt{MR-04} \\
\textbf{Decision kind} & \texttt{SELECTABLE} \\
\bottomrule
\end{tabularx}
\end{decmeta}

\subsubsection{Context}
L'adattamento della classificazione di urgenza necessita di un target di supervisione. La fonte usa la priorità P-scale clinicamente corretta come sorgente semantica per derivare il target HIGH/MEDIUM/LOW.

\subsubsection{Decision}
Per i casi di adattamento classificatorio con priorità P-scale clinicamente corretta, DermaTriage deriva il target di supervisione dalla priorità corretta secondo il mapping governato P1/P2-HIGH, P3-MEDIUM, P4-LOW.

\subsubsection{Consequences}
Il mapping P1/P2 $\rightarrow$ HIGH, P3 $\rightarrow$ MEDIUM e P4 $\rightarrow$ LOW fornisce il target di supervisione al percorso di adattamento classificatorio.

\subsubsection{FR-08 - Derivazione del target di supervisione dalla priorità P-scale corretta}\label{fr-08}

\begin{frmeta}
\begin{tabularx}{\linewidth}{L{43mm}Y}
\toprule
\textbf{ID} & \texttt{FR-08} \\
\textbf{Lifecycle} & \texttt{current} \\
\textbf{Authority} & \candidateauthority \\
\textbf{Type} & \FR \\
\textbf{Title} & Derivazione del target di supervisione dalla priorità P-scale corretta \\
\textbf{Parent Decision} & \texttt{DEC-11} \\
\textbf{Owning MR (derived)} & \texttt{MR-04} \\
\textbf{SpecializedRequirement relation} & 0 active instances \\
\bottomrule
\end{tabularx}
\end{frmeta}

\paragraph{Functional obligation.}
Derivare il target di supervisione classificatoria dalla priorità P-scale corretta dal professionista sanitario.

\begin{normbox}[title={Functional clauses / inherited \texttt{normativeClause [1..*]}}]
\begin{tabularx}{\linewidth}{L{45mm}Y}
\toprule
\textbf{Priorità clinicamente corretta} & \textbf{Target di supervisione richiesto} \\
\midrule
\texttt{P1} & DermaTriage MUST derivare \texttt{HIGH}. \\
\texttt{P2} & DermaTriage MUST derivare \texttt{HIGH}. \\
\texttt{P3} & DermaTriage MUST derivare \texttt{MEDIUM}. \\
\texttt{P4} & DermaTriage MUST derivare \texttt{LOW}. \\
\bottomrule
\end{tabularx}
\end{normbox}

\paragraph{SPO references (L2 support).}
\begin{lstlisting}
ClinicallyCorrectedOperationalTriagePriority -- maps to --> ClassifierUrgencyTarget
DermaTriage -- derive --> ClassifierUrgencyTarget
\end{lstlisting}




\subsection{DEC-17 - Gestione dei prompt con PromptManager}\label{dec-17}
\begin{decmeta}
\begin{tabularx}{\linewidth}{L{43mm}Y}
\toprule
\textbf{ID} & \texttt{DEC-17} \\
\textbf{Lifecycle} & \texttt{current} \\
\textbf{Authority} & \candidateauthority \\
\textbf{Type} & \DEC \\
\textbf{Title} & Gestione dei prompt con PromptManager \\
\textbf{Parent MacroRequirement} & \texttt{MR-04} \\
\textbf{Decision kind} & \texttt{SELECTABLE} \\
\bottomrule
\end{tabularx}
\end{decmeta}

\subsubsection{Context}
L'evoluzione dei prompt di Qwen2-VL e BioMistral richiede versioning, persistenza e operazioni di aggiornamento coerenti con il ciclo di adattamento.

\subsubsection{Decision}
DermaTriage usa \texttt{PromptManager} come componente di gestione dei prompt, con persistenza in \nolinkurl{db/prompts.json}; il componente è usato per update, reset, history e selezione della versione applicabile ai modelli Qwen2-VL e BioMistral.

\subsubsection{Consequences}
La gestione dei prompt resta distinta dal retraining classificatorio e mantiene un proprio stato di versione e persistenza.

\subsubsection{FR-26 - Versioning e persistenza dei prompt}\label{fr-26}
\begin{frmeta}
\begin{tabularx}{\linewidth}{L{43mm}Y}
\toprule
\textbf{ID} & \texttt{FR-26} \\
\textbf{Lifecycle} & \texttt{current} \\
\textbf{Authority} & \candidateauthority \\
\textbf{Type} & \FR \\
\textbf{Title} & Versioning e persistenza dei prompt \\
\textbf{Parent Decision} & \texttt{DEC-17} \\
\textbf{Owning MR (derived)} & \texttt{MR-04} \\
\textbf{SpecializedRequirement relation} & 0 active instances \\
\bottomrule
\end{tabularx}
\end{frmeta}
\paragraph{Functional obligation.} Gestire e persistere le versioni dei prompt usati dal percorso di evoluzione.
\begin{normbox}[title={Functional clauses / inherited \texttt{normativeClause [1..*]}}]
PromptManager MUST mantenere le versioni dei prompt usati da Qwen2-VL e BioMistral, MUST persisterne lo stato in \nolinkurl{db/prompts.json} e MUST supportare le operazioni documentate di update, reset e history.
\end{normbox}
\paragraph{SPO references (L2 support).}
\begin{lstlisting}
PromptManager -- manages --> PromptVersion
PromptManager -- persists to --> db/prompts.json
PromptVersion -- used by --> Qwen2-VL / BioMistral
\end{lstlisting}

\subsection{DEC-18 - Retraining incrementale di EfficientNet-B4}\label{dec-18}
\begin{decmeta}
\begin{tabularx}{\linewidth}{L{43mm}Y}
\toprule
\textbf{ID} & \texttt{DEC-18} \\
\textbf{Lifecycle} & \texttt{current} \\
\textbf{Authority} & \candidateauthority \\
\textbf{Type} & \DEC \\
\textbf{Title} & Retraining incrementale di EfficientNet-B4 \\
\textbf{Parent MacroRequirement} & \texttt{MR-04} \\
\textbf{Decision kind} & \texttt{SELECTABLE} \\
\bottomrule
\end{tabularx}
\end{decmeta}

\subsubsection{Context}
Il percorso classificatorio deve adattare EfficientNet-B4 usando correzioni clinicamente validate senza ricostruire l'intero modello da zero a ogni ciclo.

\subsubsection{Decision}
Il retraining acquisisce dai casi B4 le correzioni qualificate e i documenti applicabili, usa immagini JPEG/PNG ed esclude i PDF dall'input CNN; EfficientNet-B4 viene fine-tuned mantenendo congelati i layer precedenti e aggiornando gli ultimi due feature block e la classifier head.

\subsubsection{Consequences}
Il ciclo usa backbone learning rate 5e-6, classifier learning rate 3e-5, 10 epoch e class weights [3.0, 2.0, 1.0], con peso maggiore per HIGH.\par
Il mapping di supervisione e il trigger di attivazione restano distinti dalla strategia di fine-tuning.

\subsubsection{FR-27 - Fine-tuning degli ultimi blocchi e classifier head}\label{fr-27}
\begin{frmeta}
\begin{tabularx}{\linewidth}{L{43mm}Y}
\toprule
\textbf{ID} & \texttt{FR-27} \\
\textbf{Lifecycle} & \texttt{current} \\
\textbf{Authority} & \candidateauthority \\
\textbf{Type} & \FR \\
\textbf{Title} & Fine-tuning degli ultimi blocchi e classifier head \\
\textbf{Parent Decision} & \texttt{DEC-18} \\
\textbf{Owning MR (derived)} & \texttt{MR-04} \\
\textbf{SpecializedRequirement relation} & 0 active instances \\
\bottomrule
\end{tabularx}
\end{frmeta}
\paragraph{Functional obligation.} Eseguire il retraining incrementale di EfficientNet-B4 secondo l'algoritmo di fine-tuning selezionato.
\begin{normbox}[title={Functional clauses / inherited \texttt{normativeClause [1..*]}}]
Per ogni ciclo di retraining attivato, DermaTriage MUST fine-tune EfficientNet-B4 mantenendo congelati i layer precedenti e aggiornando gli ultimi due feature block e la classifier head; MUST usare backbone learning rate \texttt{5e-6}, classifier learning rate \texttt{3e-5}, 10 epoch e class weights \texttt{[3.0,2.0,1.0]}. I documenti B4 usati per il training MUST includere immagini JPEG/PNG applicabili e MUST escludere i PDF dall'input CNN.
\end{normbox}
\paragraph{SPO references (L2 support).}
\begin{lstlisting}
QualifiedClinicalCorrections -- supervise --> EfficientNet-B4Retraining
B4CaseImages -- provide --> CNNTrainingInput
EfficientNet-B4Retraining -- updates --> LastTwoFeatureBlocks + ClassifierHead
\end{lstlisting}

\clearpage
\section{Registro finale dei gap documentali aperti}

Questa sezione è l'unico registro dei gap della baseline corrente. Un gap indica significato non sufficientemente governato o binding incompleto; non equivale automaticamente a un nuovo Requirement.

\small
\begin{longtable}{L{50mm}L{29mm}L{64mm}}
\toprule
\textbf{Gap} & \textbf{Area} & \textbf{Significato preservato} \\
\midrule
\endhead
\texttt{Diagnostic authority boundary} & Project framing / MR-01 & L'autorità per una diagnosi clinica definitiva non è stabilita. \\
\nolinkurl{GAP-DERMA-NOIMAGE-INPUT-01} & FR-01 & Contratto completo required/optional degli input sintomatologici e semantica dei valori mancanti non specificati. \\
\nolinkurl{GAP-DERMA-PMAP-INPUT-01} & FR-02 & Semantica completa degli input urgency/confidence, inclusi valori mancanti o invalidi, non specificata. \\
\nolinkurl{GAP-DERMA-NOIMAGE-PMAP-BINDING-01} & MR-01 & Non è stabilito che l'urgenza prodotta in assenza di immagine fornisca sempre tutti gli input richiesti dal mapping P-scale. \\
\nolinkurl{GAP-DERMA-SLA-01} & MR-01 downstream & I valori 24h/48h/72h/7 giorni sono source-observed, ma trigger, outcome, owner e natura limite/target/raccomandazione non sono sufficientemente governati. \\
\nolinkurl{GAP-DERMA-PATH-01} & MR-01 downstream & Lo stato di obbligazione e il confine di autorità diagnostica di predicted\_pathology non sono stabiliti. \\
\texttt{Specialist selection semantics} & MR-02 & Vocabolario, regola di selezione, input, fallback, assignment/booking e relazione completa con P-scale non specificati. \\
\nolinkurl{GAP-DERMA-REVIEW-CONTENT-01} & FR-03 / FR-12 & Modello completo del contenuto di revisione/correzione non specificato. \\
\nolinkurl{GAP-DERMA-REVIEW-LIFECYCLE-01} & FR-03 / FR-12 & Semantica overwrite/history/retention/finality della revisione non specificata. \\
\nolinkurl{GAP-DERMA-REVIEW-DISAGREEMENT-BINDING-01} & MR-03 / MR-04 & Relazione normativa tra correzione, disaccordo e stato source-observed agrees non completamente specificata. \\
\nolinkurl{GAP-DERMA-ADAPT-COUNTING-01} & FR-04 / FR-05 & Semantica di reset, deduplication, reuse, persistence, review multiple e concurrency dei conteggi non specificata. \\
\nolinkurl{GAP-DERMA-ADAPT-THRESHOLD-AUTH-01} & FR-04 / FR-05 & Autorità di modifica e status normativo esatto dei threshold concreti non specificati. \\
\nolinkurl{GAP-DERMA-PROMPT-WINDOW-01} & FR-06 & Ordering, membership, underfill, deduplication, reuse/overlap e scope esatto della finestra non specificati. \\
\nolinkurl{GAP-DERMA-SUPERVISION-INPUT-01} & FR-08 & Comportamento per priorità corretta mancante, invalida, fuori dominio o conflittuale non specificato. \\
\nolinkurl{GAP-DERMA-DEPLOY-01} & FR-09 lifecycle & Autorità finale e automaticità del deployment dopo qualification non specificate. \\
\nolinkurl{GAP-DERMA-ACCEPT-BINDING-01} & DEC-06 / DEC-07 & Binding quantitativo completo dei criteri di accettazione, inclusi reference/evaluation population e interpretazione del 5\%, non specificato. \\
\nolinkurl{GAP-DERMA-ROLLBACK-BINDING-01} & DEC-08 / FR-10 & Unità, reference, popolazione, observation window, timing, automaticità, authorization e target esatto del rollback non sufficientemente governati. \\
\nolinkurl{GAP-DERMA-EVAL-CONSISTENCY-01} & FR-09 / DEC-07 & Il model report riporta HIGH sensitivity 90.24\% ma anche HIGH recall 81.71\% e 134/164 HIGH corretti; non è specificato se derivino da modalità di valutazione differenti. \\
\bottomrule
\end{longtable}
\normalsize

\end{document}
